% GENERATED FILE. Edit canonical lessons, then run npm run generate:books.
% canonical-source-hash: fnv1a64:c6dcc20f4f8c50ce
% canonical-lessons: FA-C325-zamin-khordan, FA-C325-pich-khordan, FA-C325-asib-didan, FA-C325-sadame-didan, FA-C325-shekastegi

\chapter{Falls and injuries}
\label{ch:fa-falls-and-injuries}
\hlchaptermodality{325}

\begin{chapteropening}
\textbf{By the end of this chapter:} \emph{I can say that someone fell over, twisted an ankle, got hurt or harmed, or broke a bone.}

Complete the last lesson of chapter 325: I can say that someone fell over, twisted an ankle, got hurt or harmed, or broke a bone.
\end{chapteropening}

\section[\texorpdfstring{\fa{زمین} \fa{خوردن} (zamin khordan)}{زمین خوردن (zamin khordan)}]{\fa{زمین} \fa{خوردن} (zamin khordan) — to fall over, to fall down}

\label{lesson:FA-C325-zamin-khordan}



\begin{quote}
Before the new one: say the Persian for to collide, then the Persian for to overturn.
\end{quote}

\subsection*{You'll want to know: \fa{زمین} \fa{خوردن}}
\textbf{\fa{زمین} \fa{خوردن}} (\emph{zamin khordan}) — \textquotedblleft{}to fall over, to fall down\textquotedblright{}.

\begin{quote}
\fa{پیرمرد} \fa{روی} \fa{یخ} \fa{زمین} \fa{خورد}.
\end{quote}

\begin{quote}
\emph{Pirmard ru-ye yakh zamin khord.}
\end{quote}

\begin{quote}
The old man fell over on the ice.
\end{quote}

Literally \textquotedblleft{}to hit the ground\textquotedblright{}. Present \emph{zamin mikhorad}.

\begin{cousinweb}
\emph{zamin} \textquotedblleft{}ground, earth\textquotedblright{} and \emph{khordan} \textquotedblleft{}to eat, to strike\textquotedblright{} are both native Persian; EtymDB takes \emph{khordan} back to Middle Persian \emph{xwar-}.
\end{cousinweb}

\begin{grammarlens}[title={What you've built}]
One more word for telling, explaining and describing at B1.
\end{grammarlens}

\subsection*{Your turn}
\begin{itemize}
\raggedright
  \item \emph{Say it:} \emph{zamin khordan}
  \item \emph{Say it:} \emph{zamin khordan}, once more
  \item \emph{Say it:} a sentence of your own with \emph{zamin khordan}
  \item \emph{Recall:} say the Persian for to collide, then the Persian for to overturn, then say \emph{zamin khordan} again
\end{itemize}

\begin{tcolorbox}[breakable,colback=teal!4,colframe=teal!35!black,title={Before you move on}]
What does \fa{زمین} \fa{خوردن} mean? (\textquotedblleft{}To fall over, to fall down\textquotedblright{}.) Say it once more.
\end{tcolorbox}
\section[\texorpdfstring{\fa{پیچ} \fa{خوردن} (pich khordan)}{پیچ خوردن (pich khordan)}]{\fa{پیچ} \fa{خوردن} (pich khordan) — to get twisted, to sprain}

\label{lesson:FA-C325-pich-khordan}



\begin{quote}
Before the new one: say the Persian for to overturn, then the Persian for to fall over.
\end{quote}

\subsection*{You'll want to know: \fa{پیچ} \fa{خوردن}}
\textbf{\fa{پیچ} \fa{خوردن}} (\emph{pich khordan}) — \textquotedblleft{}to get twisted, to sprain\textquotedblright{}.

\begin{quote}
\fa{پایم} \fa{در} \fa{کوه} \fa{پیچ} \fa{خورد}.
\end{quote}

\begin{quote}
\emph{Pâyam dar kuh pich khord.}
\end{quote}

\begin{quote}
I twisted my ankle on the mountain.
\end{quote}

The body part is the subject: \emph{dastam pich khord} \textquotedblleft{}I sprained my wrist\textquotedblright{}. Present \emph{pich mikhorad}.

\begin{cousinweb}
\emph{pich} \textquotedblleft{}a twist, a turn\textquotedblright{} is a native Persian word that EtymDB lists without a parent; \emph{khordan} \textquotedblleft{}to eat, to strike\textquotedblright{} is inherited Persian.
\end{cousinweb}

\begin{grammarlens}[title={What you've built}]
One more word for telling, explaining and describing at B1.
\end{grammarlens}

\subsection*{Your turn}
\begin{itemize}
\raggedright
  \item \emph{Say it:} \emph{pich khordan}
  \item \emph{Say it:} \emph{pich khordan}, once more
  \item \emph{Say it:} a sentence of your own with \emph{pich khordan}
  \item \emph{Recall:} say the Persian for to overturn, then the Persian for to fall over, then say \emph{pich khordan} again
\end{itemize}

\begin{tcolorbox}[breakable,colback=teal!4,colframe=teal!35!black,title={Before you move on}]
What does \fa{پیچ} \fa{خوردن} mean? (\textquotedblleft{}To get twisted, to sprain\textquotedblright{}.) Say it once more.
\end{tcolorbox}
\section[\texorpdfstring{\fa{آسیب} \fa{دیدن} (âsib didan)}{آسیب دیدن (âsib didan)}]{\fa{آسیب} \fa{دیدن} (âsib didan) — to get hurt, to be damaged}

\label{lesson:FA-C325-asib-didan}



\begin{quote}
Before the new one: say the Persian for to fall over, then the Persian for to get twisted.
\end{quote}

\subsection*{You'll want to know: \fa{آسیب} \fa{دیدن}}
\textbf{\fa{آسیب} \fa{دیدن}} (\emph{âsib didan}) — \textquotedblleft{}to get hurt, to be damaged\textquotedblright{}.

\begin{quote}
\fa{دو} \fa{نفر} \fa{در} \fa{این} \fa{حادثه} \fa{آسیب} \fa{دیدند}.
\end{quote}

\begin{quote}
\emph{Do nafar dar in hâdese âsib didand.}
\end{quote}

\begin{quote}
Two people were hurt in this accident.
\end{quote}

Literally \textquotedblleft{}to see harm\textquotedblright{}. Present stem \emph{bin-}: \emph{âsib mibinad}.

\begin{cousinweb}
\emph{âsib} \textquotedblleft{}harm, injury\textquotedblright{} and \emph{didan} \textquotedblleft{}to see\textquotedblright{} are both native Persian words; EtymDB lists no parent for either.
\end{cousinweb}

\begin{grammarlens}[title={What you've built}]
One more word for telling, explaining and describing at B1.
\end{grammarlens}

\subsection*{Your turn}
\begin{itemize}
\raggedright
  \item \emph{Say it:} \emph{âsib didan}
  \item \emph{Say it:} \emph{âsib didan}, once more
  \item \emph{Say it:} a sentence of your own with \emph{âsib didan}
  \item \emph{Recall:} say the Persian for to fall over, then the Persian for to get twisted, then say \emph{âsib didan} again
\end{itemize}

\begin{tcolorbox}[breakable,colback=teal!4,colframe=teal!35!black,title={Before you move on}]
What does \fa{آسیب} \fa{دیدن} mean? (\textquotedblleft{}To get hurt, to be damaged\textquotedblright{}.) Say it once more.
\end{tcolorbox}
\section[\texorpdfstring{\fa{صدمه} \fa{دیدن} (sadame didan)}{صدمه دیدن (sadame didan)}]{\fa{صدمه} \fa{دیدن} (sadame didan) — to be harmed, to suffer damage}

\label{lesson:FA-C325-sadame-didan}



\begin{quote}
Before the new one: say the Persian for to get twisted, then the Persian for to get hurt.
\end{quote}

\subsection*{You'll want to know: \fa{صدمه} \fa{دیدن}}
\textbf{\fa{صدمه} \fa{دیدن}} (\emph{sadame didan}) — \textquotedblleft{}to be harmed, to suffer damage\textquotedblright{}.

\begin{quote}
\fa{ماشین} \fa{در} \fa{تصادف} \fa{خیلی} \fa{صدمه} \fa{دید}.
\end{quote}

\begin{quote}
\emph{Mâshin dar tasâdof kheyli sadame did.}
\end{quote}

\begin{quote}
The car was badly damaged in the crash.
\end{quote}

\emph{Sadame zadan} is \textquotedblleft{}to harm\textquotedblright{}. Close in meaning to \emph{âsib didan}.

\begin{cousinweb}
\emph{sadame} is borrowed from Arabic \emph{ṣadma} (EtymDB lists the Arabic word), usually glossed \textquotedblleft{}a blow, a shock\textquotedblright{}. \emph{didan} \textquotedblleft{}to see\textquotedblright{} is native Persian.
\end{cousinweb}

\begin{grammarlens}[title={What you've built}]
One more word for telling, explaining and describing at B1.
\end{grammarlens}

\subsection*{Your turn}
\begin{itemize}
\raggedright
  \item \emph{Say it:} \emph{sadame didan}
  \item \emph{Say it:} \emph{sadame didan}, once more
  \item \emph{Say it:} a sentence of your own with \emph{sadame didan}
  \item \emph{Recall:} say the Persian for to get twisted, then the Persian for to get hurt, then say \emph{sadame didan} again
\end{itemize}

\begin{tcolorbox}[breakable,colback=teal!4,colframe=teal!35!black,title={Before you move on}]
What does \fa{صدمه} \fa{دیدن} mean? (\textquotedblleft{}To be harmed, to suffer damage\textquotedblright{}.) Say it once more.
\end{tcolorbox}
\section[\texorpdfstring{\fa{شکستگی} (shekastegi)}{شکستگی (shekastegi)}]{\fa{شکستگی} (shekastegi) — a fracture, a break}

\label{lesson:FA-C325-shekastegi}



\begin{quote}
Before the new one: say the Persian for to get hurt, then the Persian for to be harmed.
\end{quote}

\subsection*{You'll want to know: \fa{شکستگی}}
\textbf{\fa{شکستگی}} (\emph{shekastegi}) — \textquotedblleft{}a fracture, a break\textquotedblright{}.

\begin{quote}
\fa{دکتر} \fa{گفت} \fa{شکستگی} \fa{پایش} \fa{جدی} \fa{نیست}.
\end{quote}

\begin{quote}
\emph{Doktor goft shekastegi-ye pâyash jeddi nist.}
\end{quote}

\begin{quote}
The doctor said the break in his leg was not serious.
\end{quote}

From \emph{shekaste} \textquotedblleft{}broken\textquotedblright{}, the participle of \emph{shekastan} \textquotedblleft{}to break\textquotedblright{}.

\begin{cousinweb}
\emph{shekastegi} is built on \emph{shekaste} \textquotedblleft{}broken\textquotedblright{}, the participle of the native verb \emph{shekastan} \textquotedblleft{}to break\textquotedblright{}, with the noun suffix \emph{-gi}.
\end{cousinweb}

\begin{grammarlens}[title={What you've built}]
One more word for telling, explaining and describing at B1.
\end{grammarlens}

\subsection*{Your turn}
\begin{itemize}
\raggedright
  \item \emph{Say it:} \emph{shekastegi}
  \item \emph{Say it:} \emph{shekastegi}, once more
  \item \emph{Say it:} a sentence of your own with \emph{shekastegi}
  \item \emph{Recall:} say the Persian for to get hurt, then the Persian for to be harmed, then say \emph{shekastegi} again
\end{itemize}

\begin{tcolorbox}[breakable,colback=teal!4,colframe=teal!35!black,title={Before you move on}]
What does \fa{شکستگی} mean? (\textquotedblleft{}A fracture, a break\textquotedblright{}.) Say it once more.
\end{tcolorbox}
